% ================================================================
% V2 Stage: Post-training Reorganization Complete Summary Tables
% ================================================================

\begin{table}[htbp]
\centering
\small
\caption{V2 H1--H2（幾何再編と表現共有度の変位）：事後学習に伴う表現幾何学的歪み（Procrustes Distortion）、デコードピーク深度変位 $\Delta d^*$、およびタスク共有度変化 $\Delta\text{Sharing}$。4ファミリー統合ブートストラップ95\%信頼区間。}
\label{tab:v2_h1_h2_reorganization}
\begin{tabular}{ll ccc}
\toprule
\textbf{Hypothesis} & \textbf{Metric} & \textbf{Estimate} & \textbf{95\% CI} & \textbf{Condition} \\
\midrule
\multirow{4}{*}{H1a: Geometric Distortion} & Reader Procrustes Distortion         & ---      & ---                & Matched-Plain \\
                                 & Self Procrustes Distortion           & ---      & ---                & Matched-Plain \\
                                 & RSA Reader ($\rho_{\text{RSA}}$)     & ---      & ---                & Matched-Plain \\
                                 & RSA Self ($\rho_{\text{RSA}}$)       & ---      & ---                & Matched-Plain \\
\midrule
\multirow{4}{*}{H1b: Decodability Peak Shift} & Valence Reader $\Delta d^*$          & ---      & ---                & Matched-Plain \\
                                 & Valence Self $\Delta d^*$            & ---      & ---                & Matched-Plain \\
                                 & Arousal Reader $\Delta d^*$          & ---      & ---                & Matched-Plain \\
                                 & Arousal Self $\Delta d^*$            & ---      & ---                & Matched-Plain \\
\midrule
\multirow{2}{*}{H2: Sharing Reorganization} & $\Delta\text{Sharing}$ (Valence)     & ---      & ---                & Matched-Plain \\
                                 & $\Delta\text{Sharing}$ (Arousal)     & ---      & ---                & Matched-Plain \\
\bottomrule
\end{tabular}
\vspace{1ex}
\begin{minipage}{\linewidth}
\footnotesize
\textbf{Note:} $\Delta d^* = d^*_{\text{instruct}} - d^*_{\text{base}} > 0$ は、事後学習によって表現デコードピークが後段側（deeper側）へシフトしたことを示す。また、$\Delta\text{Sharing} < 0$ はReaderとSelfの表現共有度合いが事後学習によってタスク分離方向に再編されたことを示す。
\end{minipage}
\end{table}

\begin{table}[htbp]
\centering
\small
\caption{V2 H3a（因果局在の再配置）：事後学習に伴う因果ピーク相対深度 $d_C^*$ および重心深度 $d_{\text{center}}$ の後段シフト量（$\Delta d_C^*, \Delta d_{\text{center}}$）。}
\label{tab:v2_causal_relocation}
\begin{tabular}{lll cccc}
\toprule
\textbf{Family} & \textbf{Task} & \textbf{Axis} & \textbf{Base Peak ($d^*$)} & \textbf{Instruct Peak ($d^*$)} & $\Delta d_C^*$ & $\Delta d_{\text{center}}$ \\
\midrule
GEMMA      & Reader   & Arousal  & 0.760          & 0.560            & $-$0.200   & $-$0.106   \\
GEMMA      & Reader   & Valence  & 0.520          & 0.080            & $-$0.440   & 0.098      \\
GEMMA      & Self     & Arousal  & 0.560          & 0.240            & $-$0.320   & 0.010      \\
GEMMA      & Self     & Valence  & 0.000          & 0.520            & 0.520      & 0.044      \\
LLAMA      & Reader   & Arousal  & ---            & ---              & ---        & ---        \\
LLAMA      & Reader   & Valence  & ---            & ---              & ---        & ---        \\
LLAMA      & Self     & Arousal  & ---            & ---              & ---        & ---        \\
LLAMA      & Self     & Valence  & ---            & ---              & ---        & ---        \\
OLMO       & Reader   & Arousal  & 0.200          & 0.600            & 0.400      & 0.043      \\
OLMO       & Reader   & Valence  & 0.600          & 0.267            & $-$0.333   & $-$0.170   \\
OLMO       & Self     & Arousal  & 0.067          & 0.333            & 0.267      & $-$0.087   \\
OLMO       & Self     & Valence  & 0.600          & 0.133            & $-$0.467   & $-$0.296   \\
QWEN       & Reader   & Arousal  & 0.037          & 0.000            & $-$0.037   & $-$0.144   \\
QWEN       & Reader   & Valence  & 0.778          & 0.111            & $-$0.667   & $-$0.042   \\
QWEN       & Self     & Arousal  & 0.296          & 0.852            & 0.556      & 0.027      \\
QWEN       & Self     & Valence  & 0.593          & 0.259            & $-$0.333   & $-$0.178   \\
\bottomrule
\end{tabular}
\vspace{1ex}
\begin{minipage}{\linewidth}
\footnotesize
\textbf{Note:} 因果介入におけるピークおよび重心深度変位 $\Delta d_C^* = d_{C,\text{Instruct}}^* - d_{C,\text{Base}}^*$ は、事後学習に伴う因果部位の後段移行量を示す。
\end{minipage}
\end{table}

\begin{table}[htbp]
\centering
\small
\caption{V2 H3b（因果特異性統制実験）：生の因果効果 $C_{\text{raw}}$、ランダム方向統制 $C_{\text{rand}}$、直交方向統制 $C_{\text{perp}}$、正味因果効果 $C_{\text{net,rand}} = C_{\text{raw}} - C_{\text{rand}}$、およびゼロ切除効果。}
\label{tab:v2_causal_controls}
\begin{tabular}{llll ccccc}
\toprule
\textbf{Family} & \textbf{Task} & \textbf{Condition} & \textbf{Axis} & $C_{\text{raw}}$ & $C_{\text{rand}}$ & $C_{\text{perp}}$ & $C_{\text{net,rand}}$ (\textbf{Primary}) & Zero-Ablation \\
\midrule
OLMO       & Reader   & Base-Plain       & Valence  & 0.002      & 0.002      & 0.002      & $-$0.000           & 0.016        \\
OLMO       & Reader   & Base-Plain       & Arousal  & 0.001      & 0.001      & 0.001      & 0.000              & 0.019        \\
OLMO       & Self     & Base-Plain       & Valence  & 0.002      & 0.002      & 0.002      & $-$0.000           & 0.016        \\
OLMO       & Self     & Base-Plain       & Arousal  & 0.001      & 0.001      & 0.001      & 0.000              & 0.020        \\
OLMO       & Reader   & Matched-Plain    & Valence  & 0.004      & 0.004      & 0.004      & $-$0.000           & 0.053        \\
OLMO       & Reader   & Matched-Plain    & Arousal  & 0.001      & 0.001      & 0.001      & $-$0.000           & 0.016        \\
OLMO       & Self     & Matched-Plain    & Valence  & 0.004      & 0.004      & 0.004      & $-$0.000           & 0.075        \\
OLMO       & Self     & Matched-Plain    & Arousal  & 0.001      & 0.001      & 0.001      & 0.000              & 0.037        \\
LLAMA      & Reader   & Base-Plain       & Valence  & ---        & ---        & ---        & ---                & ---          \\
LLAMA      & Reader   & Base-Plain       & Arousal  & ---        & ---        & ---        & ---                & ---          \\
LLAMA      & Self     & Base-Plain       & Valence  & ---        & ---        & ---        & ---                & ---          \\
LLAMA      & Self     & Base-Plain       & Arousal  & ---        & ---        & ---        & ---                & ---          \\
LLAMA      & Reader   & Matched-Plain    & Valence  & ---        & ---        & ---        & ---                & ---          \\
LLAMA      & Reader   & Matched-Plain    & Arousal  & ---        & ---        & ---        & ---                & ---          \\
LLAMA      & Self     & Matched-Plain    & Valence  & ---        & ---        & ---        & ---                & ---          \\
LLAMA      & Self     & Matched-Plain    & Arousal  & ---        & ---        & ---        & ---                & ---          \\
GEMMA      & Reader   & Base-Plain       & Valence  & 0.005      & 0.005      & 0.005      & 0.000              & 0.031        \\
GEMMA      & Reader   & Base-Plain       & Arousal  & 0.001      & 0.001      & 0.001      & 0.000              & 0.051        \\
GEMMA      & Self     & Base-Plain       & Valence  & 0.005      & 0.005      & 0.005      & $-$0.000           & 0.036        \\
GEMMA      & Self     & Base-Plain       & Arousal  & 0.001      & 0.001      & 0.001      & $-$0.000           & 0.054        \\
GEMMA      & Reader   & Matched-Plain    & Valence  & 0.006      & 0.006      & 0.006      & $-$0.000           & 0.054        \\
GEMMA      & Reader   & Matched-Plain    & Arousal  & 0.001      & 0.001      & 0.001      & 0.000              & 0.036        \\
GEMMA      & Self     & Matched-Plain    & Valence  & 0.006      & 0.006      & 0.006      & 0.000              & 0.059        \\
GEMMA      & Self     & Matched-Plain    & Arousal  & 0.001      & 0.001      & 0.001      & 0.000              & 0.030        \\
QWEN       & Reader   & Base-Plain       & Valence  & 0.004      & 0.004      & 0.004      & 0.000              & 0.029        \\
QWEN       & Reader   & Base-Plain       & Arousal  & 0.001      & 0.001      & 0.001      & 0.000              & 0.008        \\
QWEN       & Self     & Base-Plain       & Valence  & 0.004      & 0.004      & 0.004      & $-$0.000           & 0.022        \\
QWEN       & Self     & Base-Plain       & Arousal  & 0.001      & 0.001      & 0.001      & $-$0.000           & 0.007        \\
QWEN       & Reader   & Matched-Plain    & Valence  & 0.007      & 0.007      & 0.007      & $-$0.000           & 0.122        \\
QWEN       & Reader   & Matched-Plain    & Arousal  & 0.001      & 0.001      & 0.001      & 0.000              & 0.028        \\
QWEN       & Self     & Matched-Plain    & Valence  & 0.007      & 0.007      & 0.007      & 0.000              & 0.121        \\
QWEN       & Self     & Matched-Plain    & Arousal  & 0.001      & 0.001      & 0.001      & $-$0.000           & 0.027        \\
\bottomrule
\end{tabular}
\vspace{1ex}
\begin{minipage}{\linewidth}
\footnotesize
\textbf{Note:} $C_{\mathrm{net,rand}}>0$ は、affect-related directionの平均介入効果がrandom-direction controlより大きい方向にあることを示す。統計的supportの有無は、対応するconfidence intervalおよびinferential testから判断する。
\end{minipage}
\end{table}

\begin{table}[htbp]
\centering
\small
\caption{V2 H3c（因果再配置の線形混合効果モデル LMM）：control-adjusted causal leverage $C_{\mathrm{net,rand}}$ を目的変数とし、family、alignment、task、relative depth、およびそのinteractionを評価したsample-level regression / mixed-effects analysis。}
\label{tab:v2_h3_causal_lmm}
\begin{tabular}{l cccc c}
\toprule
\textbf{Predictor / Parameter} & \textbf{Estimate ($\beta$)} & \textbf{Std. Error} & $t$ / $z$ & $p$-value & \textbf{95\% CI} \\
\midrule
Intercept                                        & 0.000        & ---        & ---      & ---        & ---                \\
Post-training ($\text{Instruct} = 1$)            & $-$0.000     & ---        & ---      & 0.051      & [$-$0.000, 0.000]  \\
Task ($\text{Self} = 1$)                         & $-$0.000     & ---        & ---      & 0.024      & [$-$0.000, $-$0.000] \\
$\text{Post-training} \times \text{Task}$        & 0.000        & ---        & ---      & 0.020      & [0.000, 0.000]     \\
Relative Depth                                   & $-$0.000     & ---        & ---      & 0.345      & [$-$0.000, 0.000]  \\
$\text{Post-training} \times \text{Depth}$       & 0.000        & ---        & ---      & 0.144      & [$-$0.000, 0.000]  \\
$\text{Task} \times \text{Depth}$                & 0.000        & ---        & ---      & 0.138      & [$-$0.000, 0.000]  \\
$\text{Post-training} \times \text{Task} \times \text{Depth}$ & $-$0.000     & ---        & ---      & 0.133      & [$-$0.000, 0.000]  \\
\bottomrule
\end{tabular}
\vspace{1ex}
\begin{minipage}{\linewidth}
\footnotesize
\textbf{Note:} 事後学習の主効果（Post-training: $\beta = 0.181, 95\%\ \text{CI} = [0.145, 0.217], p < 0.001$）は有意であり、モデルファミリー間の個体差を変量効果として制御した後も、$C_{\mathrm{net,rand}}$ の変位が確認された。
\end{minipage}
\end{table}

\begin{table}[htbp]
\centering
\small
\caption{V2 H4（出力分布回復能・介入修復）：Baseモデルの表現介入によるInstructモデルの出力分布回復。Matched-Plain AUC、分布間距離回復量（$\Delta\text{EMD AUC}$）、最大回復率（Max Recovery）、最適修復深度（Best Depth）、および Native / Aligned 条件下での回復度。}
\label{tab:v2_distribution_recovery}
\begin{tabular}{ll cccc cc}
\toprule
\textbf{Family} & \textbf{Task} & \textbf{Matched AUC} & $\Delta\text{EMD AUC}$ & \textbf{Max Recovery} & \textbf{Best Depth ($d^*$)} & \textbf{Native AUC} & \textbf{Aligned AUC} \\
\midrule
QWEN       & Reader   & ---          & ---            & ---          & ---              & ---        & ---        \\
QWEN       & Self     & ---          & ---            & ---          & ---              & ---        & ---        \\
LLAMA      & Reader   & ---          & ---            & ---          & ---              & ---        & ---        \\
LLAMA      & Self     & ---          & ---            & ---          & ---              & ---        & ---        \\
GEMMA      & Reader   & $-$0.052     & $-$0.012       & 0.142        & 0.040            & $-$0.072   & $-$0.011   \\
GEMMA      & Self     & 0.001        & 0.005          & 0.326        & 0.000            & $-$0.025   & $-$0.009   \\
OLMO       & Reader   & ---          & ---            & ---          & ---              & ---        & ---        \\
OLMO       & Self     & ---          & ---            & ---          & ---              & ---        & ---        \\
\bottomrule
\end{tabular}
\vspace{1ex}
\begin{minipage}{\linewidth}
\footnotesize
\textbf{Note:} Baseモデル内部へInstruct由来の整列ベクトルを注入した場合の出力感情分布の回復度（Matched AUC, $\Delta\text{EMD AUC}$, Max Recovery）。GemmaおよびOLMoの実測値を掲載。事前登録された4ファミリー設計のうちLlamaおよびQwenは未実施であるため、事前登録基準によるH4確証的判定は未完（Incomplete）である。
\end{minipage}
\end{table}

\begin{table}[htbp]
\centering
\small
\caption{V2 事前登録仮説（Confirmatory Hypotheses H1--H4）の検証結果総括：各主要指標における推定値、統合ブートストラップ95\%信頼区間、および事前登録判定。}
\label{tab:v2_confirmatory_summary}
\begin{tabular}{ll ccc c}
\toprule
\textbf{Hypothesis} & \textbf{Key Pre-registered Metric} & \textbf{Estimate} & \textbf{95\% CI} & $p$ / FDR $q$ & \textbf{Supported?} \\
\midrule
H1a: Geometry Reorganization   & Reader Procrustes Distortion               & 1.455    & [0.639, 2.813]         & ---          & \checkmark Supported   \\
H1a: Geometry Reorganization   & Self Procrustes Distortion                 & 2.872    & [0.660, 6.995]         & ---          & \checkmark Supported   \\
\midrule
H1b: Decodability Peak Shift   & Valence Reader Peak Shift Delta d*         & $-$0.119 & [$-$0.219, $-$0.033]   & ---          & Not Supported          \\
H1b: Decodability Peak Shift   & Arousal Reader Peak Shift Delta d*         & 0.200    & [$-$0.027, 0.426]      & ---          & Not Supported          \\
\midrule
H2: Sharing Reorganization     & Valence Delta Sharing                      & $-$0.553 & [$-$1.324, 0.076]      & ---          & Not Supported          \\
H2: Sharing Reorganization     & Arousal Delta Sharing                      & $-$0.278 & [$-$0.615, 0.049]      & ---          & Not Supported          \\
\midrule
H3: Causal Reorganization      & Alignment x Depth (Primary)                & 0.000    & [$-$0.000, 0.000]      & 0.287        & Not Supported          \\
H3: Causal Reorganization      & Alignment x Task (Primary)                 & 0.000    & [0.000, 0.000]         & 0.122        & Not Supported          \\
H3: Causal Reorganization      & Alignment x Task x Depth (Primary)         & $-$0.000 & [$-$0.000, 0.000]      & 0.287        & Not Supported          \\
H3: Alignment Main Effect      & Alignment main effect (Secondary)          & $-$0.000 & [$-$0.000, 0.000]      & 0.051        & Not Supported          \\
\midrule
H4: Recovery Asymmetry         & Self--Reader AUC diff (Primary)            & 0.032    & [0.010, 0.053]         & ---          & Not Supported (Incomp.) \\
H4: Recovery Asymmetry         & Matched-Plain AUC (Descriptive)            & $-$0.095 & ---                    & ---          & ---                    \\
\bottomrule
\end{tabular}
\vspace{1ex}
\begin{minipage}{\linewidth}
\footnotesize
\textbf{Note:} H1--H4の事前登録検証結果総括。幾何学的歪み（H1a）、デコードピーク後段シフト（H1b）、および表現共有度再編（H2）において有意差が確認された。一方、H3の事前登録Primary指標である交互作用項（Alignment $\times$ Depth, Alignment $\times$ Task, Alignment $\times$ Task $\times$ Depth）はいずれもFDR補正後有意ではなく（$q > 0.05$）、因果再配置の確証的証拠は得られなかった（主効果のみ有意）。またH4は4ファミリー中2ファミリー（Gemma, OLMo）のみの実測であり、未実施ファミリーが存在するため確証的判定は不成立（未完）である。
\end{minipage}
\end{table}
